\documentclass[11pt,a4paper]{article}

% ---------- packages ----------
\usepackage[utf8]{inputenc}
\usepackage[T1]{fontenc}
\usepackage{lmodern}
\usepackage{amsmath,amssymb,amsthm}
\usepackage{mathtools}
\usepackage{physics}           % \vb, \abs, \norm
\usepackage{bm}
\usepackage{geometry}
\geometry{margin=2.5cm}
\usepackage{enumitem}
\usepackage{booktabs}
\usepackage{caption}
\usepackage[dvipsnames]{xcolor}
\usepackage{hyperref}
\hypersetup{colorlinks,linkcolor=MidnightBlue,citecolor=OliveGreen,urlcolor=NavyBlue}
\usepackage{cleveref}

% ---------- theorem styles ----------
\newtheorem{theorem}{Theorem}
\newtheorem{lemma}[theorem]{Lemma}
\newtheorem{proposition}[theorem]{Proposition}
\newtheorem{corollary}[theorem]{Corollary}
\theoremstyle{definition}
\newtheorem{definition}[theorem]{Definition}
\theoremstyle{remark}
\newtheorem{remark}[theorem]{Remark}

% ---------- shortcuts ----------
\newcommand{\Sab}{S_{\mathrm{ab}}}
\newcommand{\Sinc}{S_{\mathrm{inc}}}
\newcommand{\Tzero}{T_0}
\newcommand{\nhat}{\hat{\vb{n}}}
\newcommand{\khat}{\hat{\vb{k}}}
\newcommand{\ntilde}{\tilde{n}}
\newcommand{\epstilde}{\tilde{\varepsilon}}
\newcommand{\SAR}{\ensuremath{\mathrm{SAR}}}
\newcommand{\psSAR}{\ensuremath{\mathrm{psSAR}_{10\mathrm{g}}}}
\newcommand{\rhom}{\rho_{\mathrm{m}}}
\newcommand{\ReLU}{\operatorname{ReLU}}
\DeclareMathOperator*{\argmax}{arg\,max}

% ---------- title ----------
\title{%
  Peak spatial-average SAR over 10\,g\\
  from surface absorbed power density:\\
  a closed-form reduction%
}
\author{Robin Wydaeghe}
\date{April 2026}

\begin{document}
\maketitle

\begin{abstract}
We derive the peak spatial-average specific absorption rate over a 10\,g cube
of tissue (\psSAR) from the absorbed power density $\Sab$ defined on the body
surface.  Assuming a homogeneous lossy half-space at each surface point and
exploiting the exponential decay of the internal fields, the three-dimensional
volumetric problem reduces to a two-dimensional surface operation.

Two formulations are given.  For an arbitrary (relaxed) averaging geometry, the
depth integral is performed analytically and encoded in a scalar function
$h(u) = (1-e^{-u})/u$, yielding a one-parameter optimization over connected
surface patches.  For the standard axis-aligned cube prescribed by ICNIRP~2020
and IEEE~C95.3, an energy conservation argument gives a closed-form result in
the thin-skin limit ($\alpha L \gg 1$, always satisfied above 6\,GHz):
\[
  \langle\SAR\rangle_{\mathrm{cube}}
  = \frac{\Sinc\;\Tzero}{\rhom\,L}\;
    \bigl(|k_x| + |k_y| + |k_z|\bigr),
\]
where $L = (m/\rhom)^{1/3}$ is the cube side length and
$(k_x, k_y, k_z)$ are the components of the propagation direction $\khat$.
The cube's axis-aligned geometry, the tilted body surface, and the depth
decay are all absorbed into the projected-area factor
$|k_x| + |k_y| + |k_z|$.  The result is specialized to a sphere (where
the relaxed optimization yields a single transcendental equation) and
extended to complex anatomy (ears, fingers) where the body surface folds
within a single cube.
\end{abstract}

\tableofcontents


% ======================================================================
\section{Internal fields from surface absorbed power density}
\label{sec:internal-fields}
% ======================================================================

\subsection{Setup and notation}

Consider a plane wave with propagation direction $\khat$ and incident power
density $\Sinc$ arriving at a point $\vb{r}$ on the body surface $\Sigma$.
The outward unit normal is $\nhat(\vb{r})$ and the incidence cosine is
\begin{equation}\label{eq:mu-def}
  \mu(\vb{r}) = \nhat(\vb{r}) \cdot (-\khat) = \cos\theta_i(\vb{r}).
\end{equation}
The tissue is modelled as a lossy half-space with complex relative permittivity
\begin{equation}\label{eq:eps-complex}
  \epstilde = \varepsilon_r - \frac{j\sigma}{\omega\varepsilon_0}
\end{equation}
and complex refractive index
\begin{equation}\label{eq:ntilde}
  \ntilde = \sqrt{\epstilde} = n - j\kappa, \qquad n, \kappa > 0.
\end{equation}
The normal component of the transmitted wave vector in tissue is
\begin{equation}\label{eq:xi-def}
  \xi = \sqrt{\ntilde^2 - 1 + \mu^2}, \qquad \Re(\xi) > 0.
\end{equation}
At normal incidence $(\mu = 1)$, $\xi = \ntilde$.

\subsection{Transmitted field at depth}

The transmitted electric field at surface point $\vb{r}$ and depth $z$
(measured along the inward normal) is
\begin{equation}\label{eq:Etrans}
  \vb{E}^{\mathrm{trans}}(\vb{r}, z)
  = \bigl[t_s\,\psi_s\,\hat{\vb{e}}_s + t_p\,\psi_p\,\hat{\vb{e}}_p\bigr]\,
    e^{-jk_0\khat\cdot\vb{r}}\,e^{-jk_0\xi z}\,H(\mu),
\end{equation}
where $t_s = 2\mu/(\mu + \xi)$ and $t_p = 2\ntilde\mu/(\ntilde^2\mu + \xi)$
are the Fresnel amplitude transmission coefficients, $\psi_{s,p}$ are the
polarization amplitudes, and $H(\mu)$ is the Heaviside step function.
The depth factor decomposes as
\begin{equation}\label{eq:depth-factor}
  e^{-jk_0\xi z} = e^{-\alpha z}\,e^{-j\beta z},
\end{equation}
where $\alpha = k_0\kappa_\xi > 0$ is the amplitude decay rate and
$\beta = k_0\Re(\xi)$ is the phase rate into depth, with
$\kappa_\xi = -\Im(\xi)$.

\begin{remark}[Universal decay rate]\label{rem:universal-alpha}
Because $|\ntilde| \gg 1$ for biological tissue (e.g.\ $|\ntilde| = 4.84$ for
skin at 28\,GHz), Snell's law compresses all refraction angles to within roughly
$10^\circ$ of normal. The decay rate $\alpha$ varies by at most 2.2\% across
$\theta_i \in [0^\circ, 85^\circ]$. We therefore adopt the normal-incidence
value $\alpha \equiv \alpha_0 = k_0\kappa$ throughout, introducing at most 2\%
error in the depth profile. This is the content of the monograph's universal
depth coupling approximation.
\end{remark}

\subsection{Connecting surface fields to absorbed power density}

The absorbed power density on the surface is
\begin{equation}\label{eq:Sab}
  \Sab(\vb{r}) = \Sinc\,T_{\mathrm{eff}}(\vb{r})\,\ReLU\!\bigl(\mu(\vb{r})\bigr),
\end{equation}
with $T_{\mathrm{eff}}$ the polarization-weighted Fresnel power-absorption
coefficient. By the depth-integral identity (connecting Poynting flux conservation
to the Fresnel coefficients), the surface field intensity satisfies
\begin{equation}\label{eq:E-surface}
  |\vb{E}^{\mathrm{trans}}(\vb{r}, 0)|^2
  = \frac{4\alpha}{\sigma}\,\Sab(\vb{r}).
\end{equation}
\emph{Derivation.} Integrate the ohmic loss density $\frac{\sigma}{2}|E|^2$
over all depths:
\[
  \Sab(\vb{r})
  = \frac{\sigma}{2}\int_0^\infty |\vb{E}(\vb{r},z)|^2\,\mathrm{d}z
  = \frac{\sigma}{2}\,|\vb{E}(\vb{r},0)|^2 \int_0^\infty e^{-2\alpha z}\,\mathrm{d}z
  = \frac{\sigma}{4\alpha}\,|\vb{E}(\vb{r},0)|^2.
\]
Inverting gives \cref{eq:E-surface}.

\begin{remark}[Accuracy]
\Cref{eq:E-surface} is exact for TE polarization. For TM, there is a
multiplicative correction $C_p(\theta) = 1 + \sin^2\theta/|\xi|^2$, arising
from the normal component of the electric field inside tissue. For skin at
28\,GHz, $C_p - 1 \leq 4.5$\% and is below 2\% for $\theta \leq 60^\circ$.
\end{remark}

\subsection{Full depth-resolved field intensity}

Combining \cref{eq:depth-factor,eq:E-surface}, the electric field intensity at
any depth is
\begin{equation}\label{eq:E-depth}
  \boxed{%
    |\vb{E}(\vb{r}, z)|^2
    = \frac{4\alpha}{\sigma}\,\Sab(\vb{r})\,e^{-2\alpha z}.
  }
\end{equation}
This is the bridge from the surface quantity $\Sab$ to the volumetric field.
Tissue with shorter skin depth (larger $\alpha$) has higher surface field
intensity for the same $\Sab$ because the same total absorbed power is
concentrated into a thinner layer.


% ======================================================================
\section{Volumetric SAR from internal fields}
\label{sec:volumetric-sar}
% ======================================================================

\subsection{Local SAR}

The specific absorption rate at position $(\vb{r}, z)$ inside tissue is
\begin{equation}\label{eq:SAR-def}
  \SAR(\vb{r}, z) = \frac{\sigma\,|\vb{E}(\vb{r}, z)|^2}{2\rhom},
\end{equation}
where $\rhom$ is the tissue mass density ($\mathrm{kg/m^3}$).
Substituting \cref{eq:E-depth}:
\begin{equation}\label{eq:SAR-depth}
  \boxed{%
    \SAR(\vb{r}, z) = \frac{2\alpha}{\rhom}\,\Sab(\vb{r})\,e^{-2\alpha z}.
  }
\end{equation}
The surface SAR ($z = 0$) is simply
\begin{equation}\label{eq:SAR-surface}
  \SAR(\vb{r}, 0) = \frac{2\alpha}{\rhom}\,\Sab(\vb{r}).
\end{equation}

\begin{remark}[Energy conservation check]
Integrating \cref{eq:SAR-depth} over depth and multiplying by $\rhom$ recovers
$\Sab$:
\[
  \int_0^\infty \SAR(\vb{r},z)\,\rhom\,\mathrm{d}z
  = 2\alpha\,\Sab(\vb{r})\int_0^\infty e^{-2\alpha z}\,\mathrm{d}z
  = \Sab(\vb{r}).
\]
\end{remark}

\subsection{Depth profile}

Both $|\vb{E}|^2$ and SAR decay as $e^{-2\alpha z}$. The characteristic
SAR depth is
\begin{equation}\label{eq:delta-SAR}
  \delta_{\SAR} = \frac{1}{2\alpha}.
\end{equation}
At 28\,GHz for skin: $\delta_{\SAR} \approx 0.48$\,mm. The fraction of total
absorbed power contained within the first $d$ of tissue is
$1 - e^{-2\alpha d}$~---~see \cref{tab:depth-fraction}.

\begin{table}[ht]
\centering
\caption{Fraction of absorbed power within depth $d$ at 28\,GHz (skin,
$\alpha = 1047\;\mathrm{m^{-1}}$).}
\label{tab:depth-fraction}
\begin{tabular}{@{}lcccc@{}}
\toprule
$d$ (mm) & 0.5 & 1.0 & 2.0 & 5.0 \\
\midrule
$1 - e^{-2\alpha d}$ & 0.65 & 0.88 & 0.985 & $\approx 1$ \\
\bottomrule
\end{tabular}
\end{table}


% ======================================================================
\section{Peak SAR over 10~g as a surface functional}
\label{sec:psSAR}
% ======================================================================

\subsection{Definition}

The peak spatial-average SAR over 10\,g of tissue is
\begin{equation}\label{eq:psSAR-def}
  \psSAR
  = \sup_\Omega\;\frac{1}{m_\Omega}
    \int_\Omega \frac{\sigma}{2}\,|\vb{E}|^2\,\mathrm{d}V,
\end{equation}
where $m_\Omega = \rhom\operatorname{Vol}(\Omega) = m = 10\;\mathrm{g} = 0.01\;\mathrm{kg}$.
Under ICNIRP~1998, $\Omega$ was any contiguous tissue region; under
ICNIRP~2020, $\Omega$ is restricted to a cube \cite{ICNIRP2020}.

The shape of $\Omega$ matters.  We treat two cases:
\begin{enumerate}[nosep]
\item \textbf{Surface-conforming slab} (\cref{sec:slab}): $\Omega$ is a
  connected surface patch of area $A$ extruded to depth $d = m/(\rhom A)$.
  The patch area is a free variable, yielding a one-parameter optimization.
  This is a relaxation of the standard and gives an upper bound on the
  cube-based result.
\item \textbf{Standard cube} (\cref{sec:cube}): $\Omega$ is a cube of
  side $L = (m/\rhom)^{1/3}$, as prescribed by IEEE~C95.3 and
  ICNIRP~2020.  The geometry is fixed; only the position is optimized.
\end{enumerate}

\subsection{Case~I: surface-conforming slab (relaxed geometry)}
\label{sec:slab}

At millimeter-wave frequencies the skin depth (0.3--0.5\,mm) is much smaller
than the body's minimum radius of curvature (${\sim}\,2$\,mm for an ear edge,
${\sim}\,40$\,mm for an arm). The optimal 10\,g averaging volume is therefore a
thin, surface-conforming slab: a connected surface patch $A_\Omega \subset \Sigma$
with depth extent $d_\Omega$, subject to the mass constraint
\begin{equation}\label{eq:mass-constraint}
  \rhom\,A_\Omega\,d_\Omega = m.
\end{equation}
Curvature corrections to the slab volume are $\mathcal{O}(d/R)$, where $R$ is
the local curvature radius, and amount to at most ${\sim}\,1$\% for the torso
and limbs.

Using \cref{eq:SAR-depth}, the mass-averaged SAR over the slab
$\Omega = A_\Omega \times [0, d_\Omega]$ is
\begin{align}
  \frac{1}{m}\int_\Omega \SAR\,\rhom\,\mathrm{d}V
  &= \frac{1}{m}\int_{A_\Omega}
     \frac{2\alpha}{\rhom}\,\Sab(\vb{r})\,\rhom
     \int_0^{d_\Omega} e^{-2\alpha z}\,\mathrm{d}z\;\mathrm{d}A \notag\\
  &= \frac{1}{m}\int_{A_\Omega}\Sab(\vb{r})\,
     \bigl[1 - e^{-2\alpha d_\Omega}\bigr]\,\mathrm{d}A.
     \label{eq:slab-SAR}
\end{align}

\subsection{The depth efficiency function}

\begin{definition}[Depth efficiency function]\label{def:h}
\begin{equation}\label{eq:h-def}
  h(u) = \frac{1 - e^{-u}}{u}, \qquad u > 0,
\end{equation}
with $\lim_{u \to 0^+} h(u) = 1$.
\end{definition}

The function $h$ is strictly decreasing, concave, and satisfies $h(u) \in (0, 1]$
for $u \geq 0$. It captures the fraction of surface SAR that survives
mass-averaging to a depth $d = u/(2\alpha)$.

Introduce the dimensionless depth parameter
\begin{equation}\label{eq:u-def}
  u = 2\alpha\,d_\Omega = \frac{2\alpha\,m}{\rhom\,A_\Omega},
\end{equation}
and the area-averaged absorbed power density over the patch,
\begin{equation}\label{eq:Sab-avg}
  \langle \Sab \rangle_{A_\Omega}
  = \frac{1}{A_\Omega}\int_{A_\Omega}\Sab(\vb{r})\,\mathrm{d}A.
\end{equation}
Then \cref{eq:slab-SAR} becomes
\begin{equation}\label{eq:slab-SAR-h}
  \frac{1}{m}\int_\Omega \SAR\,\rhom\,\mathrm{d}V
  = \frac{2\alpha}{\rhom}\;\langle \Sab \rangle_{A_\Omega}\;
    h\!\Bigl(\frac{2\alpha\,m}{\rhom\,A_\Omega}\Bigr).
\end{equation}

\subsection{Main result (slab)}

\begin{theorem}[psSAR from surface APD, relaxed geometry]\label{thm:main}
Under the homogeneous half-space and universal decay rate approximations
(\cref{rem:universal-alpha}),
\begin{equation}\label{eq:main}
  \boxed{%
    \psSAR
    = \frac{2\alpha}{\rhom}\;
      \sup_{\substack{A \subset \Sigma \\ A\;\text{connected},\; A > 0}}
      \left[
        \langle \Sab \rangle_A\;
        h\!\Bigl(\frac{2\alpha\,m}{\rhom\,A}\Bigr)
      \right],
  }
\end{equation}
where $\alpha = k_0\kappa$ is the normal-incidence field decay rate in tissue,
$m = 0.01\;\mathrm{kg}$, and the supremum is over connected patches on
$\Sigma$.
\end{theorem}

\begin{remark}[Physical interpretation]
The expression inside the supremum balances two competing effects:
\begin{itemize}[nosep]
\item \textbf{Small patches} (large $u$, small $h$): the area-averaged $\Sab$
  is large (it captures the hotspot), but the slab extends deep into tissue
  where SAR has decayed exponentially, reducing the mass-averaged value.
\item \textbf{Large patches} (small $u$, $h \approx 1$): the thin slab captures
  mostly high-SAR tissue near the surface, but $\langle \Sab \rangle$ is
  diluted by low-valued regions far from the hotspot.
\end{itemize}
The optimal patch area balances hotspot capture against depth dilution.
\end{remark}


% ----------------------------------------------------------------------
\subsection{Case~II: standard cube (ICNIRP/IEEE)}
\label{sec:cube}
% ----------------------------------------------------------------------

ICNIRP~1998 defined \psSAR{} using a ``10\,g contiguous tissue region'' of
unspecified shape.  ICNIRP~2020 replaced this with a \emph{cube}, explicitly
``to provide a better approximation of temperature rise''
\cite{ICNIRP2020,ICNIRPdiff}.  The $4\;\mathrm{cm}^2$ surface averaging area
was chosen to match the face of this cube, ensuring a consistent transition at
6\,GHz \cite{ICNIRPdiff}.

The computational procedure is standardized in IEC/IEEE~62704-1
(FDTD, \cite{62704-1}) and IEC/IEEE~62704-4 (FEM, \cite{62704-4}), with the
averaging volume construction specified in IEEE~C95.3-2021 Annex~E
\cite{C953}. The algorithm is:

\begin{enumerate}[nosep]
\item The computational domain is discretized into a voxel grid
  \textbf{aligned with the global Cartesian coordinate system}.
\item At each tissue voxel (seed point), a cube is \textbf{expanded evenly in
  the six axis-aligned directions} until the enclosed tissue mass is within
  $0.0001$\% of the prescribed mass $m$.
\item Validity constraints: at most 10\% of the cube volume may be background
  (air), and none of the six faces may be entirely filled with background.
\item The peak over all valid seed positions gives \psSAR.
\end{enumerate}

\noindent
Crucially, the cube is \textbf{axis-aligned}: it does not rotate to conform to
the body surface.  The phantom simply stands in the Cartesian grid and the cube
slides over voxel positions.  The cube has fixed side length
\begin{equation}\label{eq:cube-L}
  L = \Bigl(\frac{m}{\rhom}\Bigr)^{\!1/3}
  = (10^{-5}\;\mathrm{m}^3)^{1/3}
  = 21.5\;\mathrm{mm}
  \qquad (\rhom = 1000\;\mathrm{kg/m^3}).
\end{equation}
Under the half-space model, SAR decays monotonically with depth, so the optimal
cube placement has the body surface near one face.  But the cube is
axis-aligned while the body surface is not: the outward normal $\nhat$ makes
some angle $\gamma$ with the nearest cube axis.  The surface therefore cuts
through the cube as a tilted plane, and one cannot simply set $z = 0$ at the
skin.  We first treat the face-flush special case ($\gamma = 0$) as a warm-up,
then derive the general result.

\paragraph{Warm-up: face-flush cube ($\gamma = 0$).}
When the surface is parallel to a cube face, the depth integral along the
inward normal runs from $z = 0$ to $z = L$:
\begin{equation}\label{eq:cube-depth}
  \int_0^{L} \SAR(\vb{r},z)\,\rhom\,\mathrm{d}z
  = \Sab(\vb{r})\,\bigl[1 - e^{-2\alpha L}\bigr].
\end{equation}
The surface footprint is $L \times L$, the mass is $\rhom L^3$, and the
mass-averaged SAR is $(2\alpha/\rhom)\,h(2\alpha L)\,\langle\Sab\rangle_{L^2}$.
We now generalize this.

\paragraph{Tilted surface geometry.}
Orient the cube as $\mathcal{C} = [0,L]^3$ in the global $(x,y,z)$ frame.
Place $\nhat$ in the $xz$-plane without loss of generality:
$\nhat = \sin\gamma\,\hat{\vb{x}} + \cos\gamma\,\hat{\vb{z}}$,
where $\gamma \in [0, \pi/2)$ is the angle between the outward normal and
$\hat{\vb{z}}$.  The body surface $\Sigma$ is locally the plane
$\nhat \cdot \vb{r} = D$, with tissue occupying $\nhat \cdot \vb{r} < D$.
The perpendicular depth of a tissue point $\vb{r}$ below $\Sigma$ is
\begin{equation}\label{eq:depth-tilted}
  d(\vb{r}) = D - x\sin\gamma - z\cos\gamma.
\end{equation}
Using \cref{eq:SAR-depth}, the absorbed power in the tissue portion of the cube
is
\begin{equation}\label{eq:cube-3d-integral}
  \int_{\mathcal{C}\,\cap\,\text{tissue}}\!\!\SAR\,\rhom\,\mathrm{d}V
  = 2\alpha\,\Sab\;\underbrace{%
    L\!\int_0^{D}\!e^{-2\alpha(D-\zeta)}\,W(\zeta,\gamma)\,\mathrm{d}\zeta
  }_{I(\gamma,D)},
\end{equation}
where $\zeta = \nhat \cdot \vb{r}$ is the projection of $\vb{r}$ onto the
surface normal (Jacobian is unity under the orthogonal rotation
$(\zeta,\eta) \leftrightarrow (x,z)$), and $W(\zeta,\gamma)$ is the chord
length of the cube's $xz$-cross-section at level~$\zeta$.  For
$\gamma \in (0, \pi/4]$:
\begin{equation}\label{eq:W-piecewise}
  W(\zeta,\gamma) =
  \begin{dcases}
    \dfrac{\zeta}{\sin\gamma\cos\gamma}
      & 0 \le \zeta \le L\sin\gamma, \\[6pt]
    \dfrac{L}{\cos\gamma}
      & L\sin\gamma \le \zeta \le L\cos\gamma, \\[6pt]
    \dfrac{L(\sin\gamma + \cos\gamma) - \zeta}{\sin\gamma\cos\gamma}
      & L\cos\gamma \le \zeta \le L(\sin\gamma + \cos\gamma).
  \end{dcases}
\end{equation}
The middle (constant) segment is the region where the tilted plane cuts fully
across the square; its width is $L/\cos\gamma > L$.

\paragraph{Thin-skin limit ($\alpha L \gg 1$).}
At millimeter-wave frequencies, $\alpha L \approx 22$--$45$, so the integrand
$e^{-2\alpha(D-\zeta)}$ is sharply peaked at $\zeta = D$ (near the surface).
Replacing $W$ by its value at $\zeta = D$ and extending the integral to
$-\infty$:
\begin{equation}\label{eq:thin-skin-cube}
  I(\gamma, D) \approx L\,W(D,\gamma)\;\frac{1}{2\alpha}
  = \frac{L}{2\alpha}\;\frac{L}{\cos\gamma}
  = \frac{L^2}{2\alpha\cos\gamma},
\end{equation}
valid when the surface plane intersects the middle segment of
\cref{eq:W-piecewise} (i.e.\ $\gamma \lesssim 45^\circ$, which covers all but
the most oblique body orientations).

The area of the body surface intercepted by the cube is
\begin{equation}\label{eq:Agamma}
  A_\gamma = \frac{L^2}{\cos\gamma},
\end{equation}
which is the area of a tilted plane cutting through a square prism of
cross-section $L \times L$.  The face-flush case $A_0 = L^2$ is recovered at
$\gamma = 0$.

In this limit the 3D integral collapses to a surface integral.  Using
$2\alpha e^{-2\alpha d} \to \delta(d)$ as $\alpha \to \infty$:
\begin{equation}\label{eq:collapse}
  \int_{\mathcal{C}\,\cap\,\text{tissue}}\!\!\SAR\,\rhom\,\mathrm{d}V
  \;\approx\;
  \int_{\Sigma\,\cap\,\mathcal{C}} \Sab(\vb{r})\,\mathrm{d}A.
\end{equation}
All that survives is the total absorbed power flowing through the surface patch
inside the cube.

\paragraph{Mass-averaged SAR for general $\gamma$.}
Let $\langle\Sab\rangle_{A_\gamma}$ denote $\Sab$ averaged over the tilted
surface patch $\Sigma \cap \mathcal{C}$ (area $A_\gamma$).
With the tissue mass $m_{\text{tissue}} \approx \rhom L^3$ (when the air
fraction is small):
\begin{equation}\label{eq:cube-SAR-gamma}
  \frac{1}{m}\int_{\mathcal{C}\cap\text{tissue}}
    \SAR\,\rhom\,\mathrm{d}V
  \;\approx\;
  \frac{\langle\Sab\rangle_{A_\gamma}\;A_\gamma}{\rhom\,L^3}
  = \frac{\langle\Sab\rangle_{A_\gamma}}{\rhom\,L\cos\gamma}.
\end{equation}

\begin{theorem}[psSAR from surface APD, standard cube]\label{thm:cube}
In the thin-skin limit ($\alpha L \gg 1$), the \psSAR{} over axis-aligned
cubes of side $L = (m/\rhom)^{1/3}$ is
\begin{equation}\label{eq:cube-main}
  \boxed{%
    \psSAR^{\,\mathrm{cube}}
    = \frac{1}{\rhom\,L}\;
      \max_{\vb{r}_0 \in \Sigma}\;
      \frac{\langle \Sab \rangle_{A_\gamma(\vb{r}_0)}}{\cos\gamma(\vb{r}_0)},
  }
\end{equation}
where $\gamma(\vb{r}_0)$ is the angle between the surface normal at
$\vb{r}_0$ and the nearest cube axis, and
$\langle\Sab\rangle_{A_\gamma}$ is $\Sab$ averaged over the tilted surface
patch (area $L^2/\cos\gamma$) intercepted by the cube.
\end{theorem}

\begin{remark}[Geometric cancellation for aligned incidence]
For a plane wave with $\khat = -\hat{\vb{z}}$ on a convex body,
$\Sab = \Sinc\Tzero\mu$ where $\mu = \cos\theta_i$.  If $\khat$ is aligned
with a cube axis, then $\theta_i = \gamma$ at every surface point, so
$\Sab = \Sinc\Tzero\cos\gamma$.  The ratio
$\Sab / \cos\gamma = \Sinc\Tzero$ is then constant: the larger footprint at
oblique incidence exactly compensates the lower $\Sab$.  The mass-averaged SAR
$\Sinc\Tzero/(\rhom L)$ is independent of where on the body the cube is
placed.  This cancellation breaks when $\khat$ is not aligned with a cube axis,
or when Fresnel angle-dependence beyond the $\Tzero$ approximation is included.
\end{remark}

\begin{remark}[Simplification over the slab]
Unlike \cref{thm:main}, there is no optimization over patch size.
The only spatial operation is a sliding average of $\Sab$ over the
surface patch intercepted by the cube, corrected by $1/\cos\gamma$, followed
by taking its maximum.
\end{remark}

\paragraph{Numerical evaluation at 28\,GHz (skin).}
\begin{alignat}{3}
  L &= 21.5\;\mathrm{mm}, &\qquad
  1/(\rhom L) &= 1/(1000 \times 0.0215) &&= 0.047\;\mathrm{m^2/kg}, \notag\\
  \gamma &= 0 \text{ (face-flush):} &\qquad
  \psSAR^{\,\mathrm{cube}} &\approx 0.047\,\Sab^{\mathrm{peak}} &&, \notag\\
  \gamma &= 30^\circ\text{:} &\qquad
  \psSAR^{\,\mathrm{cube}} &\approx 0.054\,\langle\Sab\rangle_{A_{30}} &&.
  \label{eq:cube-numerics}
\end{alignat}
For the ICNIRP limit $\Sab = 20\;\mathrm{W/m^2}$ (uniform, face-flush):
$\psSAR^{\,\mathrm{cube}} \approx 0.93\;\mathrm{W/kg}$.

\paragraph{Energy conservation shortcut.}
The tilted-surface geometry of \cref{eq:cube-3d-integral}--\cref{eq:cube-SAR-gamma}
can be bypassed entirely by an energy conservation argument.  In the
thin-skin limit, all power transmitted through the body surface is absorbed
within the first few skin depths.  The total absorbed power inside a cube
$\mathcal{C}$ therefore equals the net inward Poynting flux through its
air-facing faces.

For a single plane wave with power density $\Sinc$ propagating along
$-\khat$, the power intercepted by the cube is
\begin{equation}\label{eq:Pabs-cube}
  P_{\mathrm{abs}} = \Sinc\;\Tzero\;A_\perp(\khat),
\end{equation}
where $A_\perp(\khat)$ is the projected area of the cube perpendicular to
$\khat$ (the ``shadow'' the cube casts on the wavefront).  For an axis-aligned
cube of side $L$, the projected area in direction
$\khat = (k_x,\,k_y,\,k_z)$ is
\begin{equation}\label{eq:Aproj}
  A_\perp(\khat) = L^2\bigl(|k_x| + |k_y| + |k_z|\bigr).
\end{equation}
When $\khat$ is along a cube axis, $A_\perp = L^2$.  At the worst case
($\khat$ along the space diagonal), $A_\perp = \sqrt{3}\,L^2 \approx 1.73\,L^2$.

The mass-averaged SAR follows immediately:
\begin{equation}\label{eq:SAR-energy}
  \boxed{%
  \langle\SAR\rangle_{\mathrm{cube}}
  = \frac{P_{\mathrm{abs}}}{m}
  = \frac{\Sinc\;\Tzero}{\rhom\,L}\;
    \bigl(|k_x| + |k_y| + |k_z|\bigr).
  }
\end{equation}
No surface tilt $\gamma$, no chord width $W(\zeta)$, no depth integral.  The
cube's axis-aligned geometry is fully encoded in the projected-area factor
$|k_x| + |k_y| + |k_z|$.  For axis-aligned incidence this reduces to
$\Sinc\Tzero/(\rhom L)$, recovering \cref{eq:cube-SAR-gamma} at $\gamma = 0$.
The geometric cancellation noted in the Remark above is now obvious: the
$\cos\gamma$ in $\Sab$ is exactly the projection factor for axis-aligned
incidence.

\Cref{eq:SAR-energy} assumes $\Tzero$ is angle-independent (the
pseudo-Brewster approximation) and that the body fills the cube's cross-section
(valid for large body parts: torso, limbs).  For small or protruding anatomy,
the shadow area may be smaller than $A_\perp$, and the full integral
$\int_{\Sigma \cap \mathcal{C}} \Sab\,\mathrm{d}A$ is needed.

\paragraph{Complex body geometry: ears, fingers, and folds.}
The preceding analysis assumed a locally planar body surface: a single,
smooth tissue--air interface crossing the cube.  Anatomy like the ear pinna
violates this completely.  An ear is roughly $60 \times 30 \times 15$\,mm;
a 21.5\,mm cube can easily enclose folds where the surface doubles back on
itself, creating tissue--air--tissue layering within a single cube.

In the thin-skin limit, the total absorbed power inside the cube remains
\begin{equation}\label{eq:Pabs-general}
  P_{\mathrm{abs}} = \int_{\Sigma\,\cap\,\mathcal{C}} \Sab(\vb{r})\,\mathrm{d}A,
\end{equation}
where the integral is over \emph{all} illuminated body-surface patches inside
$\mathcal{C}$, regardless of folding.  The mass-averaged SAR is
\begin{equation}\label{eq:SAR-general}
  \langle\SAR\rangle_{\mathrm{cube}}
  = \frac{1}{m_{\mathrm{tissue}}}
    \int_{\Sigma\,\cap\,\mathcal{C}} \Sab\,\mathrm{d}A,
  \qquad
  m_{\mathrm{tissue}} = \rhom\,V_{\mathrm{tissue}}.
\end{equation}
Two competing effects arise for folded anatomy:
\begin{itemize}[nosep]
\item \textbf{Higher surface-to-volume ratio.}  A folded structure packs more
  absorbing surface into a given cube, increasing $P_{\mathrm{abs}}$.
\item \textbf{Lower tissue mass.}  Air gaps between folds reduce
  $m_{\mathrm{tissue}}$ below $\rhom L^3$.
\end{itemize}
Both effects push $\langle\SAR\rangle$ upward, making folded anatomy a
potential hotspot.  However, there is a counteracting effect:
\textbf{self-shadowing}.  In the thin-skin limit, the outermost tissue layer
absorbs essentially all incident power; inner fold surfaces receive only the
(negligible) transmitted field.  The effective illuminated area is therefore
bounded by the cube's projected area $A_\perp(\khat)$, not the total surface
area of all folds.

For the ear specifically, the worst case occurs when the wave enters
perpendicular to the pinna, illuminating a thin cartilage sheet (thickness
${\sim}\,1$--$2$\,mm, comparable to the skin depth).  In this regime the
thin-skin approximation itself weakens: the transmitted field may partially
traverse the cartilage and illuminate the opposite surface.  A full treatment
requires accounting for finite-thickness transmission, which goes beyond the
half-space model.

\paragraph{Relationship between the two cases.}
The slab (\cref{thm:main}) relaxes the shape constraint: it allows any
aspect ratio $A/d^2$, not just $A/d^2 = 1$.  Hence the slab result is an
upper bound on the standard cube result:
\begin{equation}\label{eq:slab-geq-cube}
  \psSAR^{\,\mathrm{slab}} \geq \psSAR^{\,\mathrm{cube}}.
\end{equation}
Equality holds when the optimal slab happens to have $A = L^2$ and $d = L$.
In practice, the slab optimum favours $A \gg L^2$ (thin, wide slabs),
so the gap can be substantial: the slab captures more of the high-SAR
surface layer, while the cube wastes depth on tissue with negligible SAR.

\subsection{SAR transfer kernel}

Define the \textbf{SAR transfer kernel}
\begin{equation}\label{eq:kernel}
  \mathcal{K}_m(\vb{r}; A)
  = \frac{2\alpha}{\rhom}\;
    h\!\Bigl(\frac{2\alpha\,m}{\rhom\,A}\Bigr)\;
    \frac{\mathbf{1}_{\vb{r} \in A}}{A},
\end{equation}
so that $\psSAR = \sup_A \int_\Sigma \mathcal{K}_m(\vb{r};A)\,\Sab(\vb{r})\,\mathrm{d}A$.
This makes explicit that $\psSAR$ is a supremum over a family of smoothing
kernels applied to $\Sab$.

\subsection{Discrete form}

For a triangulated mesh with $M$ triangles (areas $a_j$, centroids $\vb{c}_j$),
\begin{equation}\label{eq:discrete}
  \psSAR
  = \frac{2\alpha}{\rhom}\;
    \max_{\substack{J \subset \{1,\ldots,M\} \\ J\;\text{connected}}}
    \left[
      \frac{\sum_{j \in J} \Sab^{(j)}\,a_j}{\sum_{j \in J} a_j}\;
      h\!\Bigl(\frac{2\alpha\,m}{\rhom\sum_{j \in J}a_j}\Bigr)
    \right].
\end{equation}

\subsection{Comparison with ICNIRP 4~cm\texorpdfstring{$^2$}{2} averaging}

The existing surface-averaging matrix $G$ from the AEGIS framework implements
ICNIRP's $4\;\mathrm{cm}^2$ averaging. For a fixed $4\;\mathrm{cm}^2$ patch
($A = 4 \times 10^{-4}\;\mathrm{m}^2$):
\[
  d = \frac{m}{\rhom\,A}
    = \frac{0.01}{1000 \times 4 \times 10^{-4}}
    = 25\;\mathrm{mm},
  \qquad
  u = 2 \times 1047 \times 0.025 = 52.4,
  \qquad
  h(52.4) = 0.019.
\]
This gives $\psSAR(4\;\mathrm{cm}^2) = 2.09 \times 0.019 \times
\langle\Sab\rangle_{4\mathrm{cm}^2} = 0.040\,\langle\Sab\rangle_{4\mathrm{cm}^2}$.
For the ICNIRP limit $\Sab = 20\;\mathrm{W/m^2}$, we obtain
$\psSAR \approx 0.80\;\mathrm{W/kg}$, comfortably below the sub-6\,GHz basic
restriction of $2\;\mathrm{W/kg}$.  ICNIRP's choice of averaging area is
conservative: $4\;\mathrm{cm}^2$ yields a 25\,mm-deep slab in which almost all
tissue contributes zero SAR.


% ======================================================================
\section{Ellipsoidal body model}
\label{sec:ellipsoid}
% ======================================================================

\subsection{Parametrization}

Consider an ellipsoid with semi-axes $a$, $b$, $c$:
\begin{equation}\label{eq:ellipsoid}
  \frac{x^2}{a^2} + \frac{y^2}{b^2} + \frac{z_{\mathrm{coord}}^2}{c^2} = 1.
\end{equation}
Parametrize by generalized spherical coordinates $(\theta, \phi)$:
\begin{equation}\label{eq:ellipsoid-param}
  \vb{r}(\theta,\phi)
  = \bigl(a\sin\theta\cos\phi,\;b\sin\theta\sin\phi,\;c\cos\theta\bigr).
\end{equation}
The outward normal direction (unnormalized) is
\begin{equation}\label{eq:N-unnorm}
  \vb{N}(\theta,\phi)
  \propto \Bigl(\frac{\sin\theta\cos\phi}{a},\;
  \frac{\sin\theta\sin\phi}{b},\;\frac{\cos\theta}{c}\Bigr),
\end{equation}
with magnitude
\begin{equation}\label{eq:N-mag}
  |\vb{N}|
  = \sqrt{\sin^2\!\theta\Bigl(\frac{\cos^2\!\phi}{a^2}
    + \frac{\sin^2\!\phi}{b^2}\Bigr) + \frac{\cos^2\!\theta}{c^2}}.
\end{equation}
The surface area element is $\mathrm{d}A = abc\,|\vb{N}|\,\sin\theta\,\mathrm{d}\theta\,\mathrm{d}\phi$.

\subsection{Plane wave along a principal axis}

Take $\khat = -\hat{\vb{z}}$ (wave illuminating the $z > 0$ hemisphere).
The incidence cosine is
\begin{equation}\label{eq:mu-ellipsoid}
  \mu(\theta,\phi)
  = \nhat \cdot \hat{\vb{z}}
  = \frac{\cos\theta}{c\,|\vb{N}(\theta,\phi)|}.
\end{equation}
For a sphere ($a = b = c = R$): $|\vb{N}| = 1/R$, recovering $\mu = \cos\theta$.
The absorbed power density under the geometric approximation is
\begin{equation}\label{eq:Sab-ellipsoid}
  \Sab(\theta,\phi) = \Sinc\,\Tzero\,\ReLU\!\bigl(\mu(\theta,\phi)\bigr).
\end{equation}

\subsection{Spherical cap optimization}
\label{sec:sphere}

For a sphere of radius $R$, consider a spherical cap of half-angle $\theta_0$
centered at the illumination pole. The cap area is
\begin{equation}\label{eq:cap-area}
  A_{\mathrm{cap}} = 2\pi R^2(1 - \cos\theta_0).
\end{equation}
The area-integrated $\Sab$ over the cap is
\begin{align}
  \int_{\mathrm{cap}} \Sab\,\mathrm{d}A
  &= \Sinc\Tzero\int_0^{\theta_0}\!\int_0^{2\pi}
     \cos\theta\;R^2\sin\theta\,\mathrm{d}\phi\,\mathrm{d}\theta \notag\\
  &= \pi R^2\Sinc\Tzero\sin^2\!\theta_0,
\end{align}
so the area-averaged $\Sab$ is
\begin{equation}\label{eq:Sab-cap}
  \langle\Sab\rangle_{\mathrm{cap}}
  = \Sinc\Tzero\;\frac{\sin^2\!\theta_0}{2(1 - \cos\theta_0)}
  = \frac{\Sinc\Tzero}{2}\,(1 + \cos\theta_0),
\end{equation}
using the identity $\sin^2\!\theta = (1-\cos\theta)(1+\cos\theta)$. As
$\theta_0 \to 0$, $\langle\Sab\rangle \to \Sinc\Tzero$ (the peak value).
At $\theta_0 = \pi/2$, $\langle\Sab\rangle = \Sinc\Tzero/2$.

The dimensionless depth parameter is
\begin{equation}\label{eq:u-cap}
  u(\theta_0) = \frac{2\alpha\,m}{\rhom\,2\pi R^2(1-\cos\theta_0)},
\end{equation}
and the objective function is
\begin{equation}\label{eq:f-sphere}
  f(\theta_0)
  = \frac{\alpha\Sinc\Tzero}{\rhom}\,(1 + \cos\theta_0)\;
    h\!\bigl(u(\theta_0)\bigr).
\end{equation}

\paragraph{Single-parameter form.}
Substituting $v = \cos\theta_0 \in [0,1]$ and defining the dimensionless
parameter
\begin{equation}\label{eq:lambda-def}
  \lambda = \frac{\alpha\,m}{\rhom\,\pi R^2},
\end{equation}
the objective becomes
\begin{equation}\label{eq:f-v}
  f(v) = \frac{\alpha\Sinc\Tzero}{\rhom}\;(1+v)\;
  h\!\Bigl(\frac{\lambda}{1-v}\Bigr).
\end{equation}
The optimal $v^*$ satisfies $f'(v^*) = 0$, a transcendental equation coupling
$v$ and $\lambda$ that must be solved numerically. The result depends on the
\emph{single} dimensionless parameter $\lambda$, which is the ratio of the SAR
depth scale $1/(2\alpha)$ to the geometric scale $m/(\rhom\pi R^2)$.

\paragraph{Numerical example.}
Take $R = 15$\,cm, $f = 28$\,GHz, $\Sinc = 10\;\mathrm{W/m^2}$,
$\Tzero = 0.539$, $\alpha = 1047\;\mathrm{m^{-1}}$, $\rhom = 1000\;\mathrm{kg/m^3}$:
\begin{itemize}[nosep]
\item $\lambda = 1047 \times 0.01 / (1000\,\pi \times 0.0225) = 0.148$.
\item Optimal cap: $\theta_0 \approx 49^\circ$,
      $A_{\mathrm{cap}} \approx 479\;\mathrm{cm}^2$,
      $d \approx 0.21$\,mm.
\item $\psSAR \approx 7.59\;\mathrm{W/kg}$.
\end{itemize}
The optimal slab is remarkably thin ($0.21$\,mm, comparable to the skin depth)
and spans a large surface area ($479\;\mathrm{cm}^2$, roughly the size of a
palm), reflecting the exponential concentration of SAR near the surface.

\subsection{Prolate spheroid}

For a prolate spheroid ($a = b < c$, a body-like shape), the azimuthal
symmetry eliminates $\phi$ and the optimization reduces to a single integral
in~$\theta$. The area element simplifies to
$\mathrm{d}A = 2\pi a\sin\theta\sqrt{a^2\cos^2\!\theta + c^2\sin^2\!\theta}\;\mathrm{d}\theta$.
The resulting one-dimensional optimization has the same structure as the sphere
case but with a modified area-versus-angle relationship that depends on the
eccentricity $e = \sqrt{1 - a^2/c^2}$.

For the triaxial case ($a \neq b \neq c$), the integrals involve elliptic
functions and the cap is no longer axially symmetric. The optimization over
two-dimensional connected patches on the triaxial ellipsoid does not admit a
closed form and requires numerical computation.


% ======================================================================
\section{Bounds and closed-form limits}
\label{sec:bounds}
% ======================================================================

\subsection{Upper bound}

Since $h(u) \leq 1$ for all $u \geq 0$ and $\langle\Sab\rangle_A \leq \max\Sab$
for any patch~$A$,
\begin{equation}\label{eq:upper-bound}
  \boxed{%
    \psSAR \leq \frac{2\alpha}{\rhom}\,\max_{\vb{r}\in\Sigma}\Sab(\vb{r})
    \leq \frac{2\alpha}{\rhom}\,\Sinc\,\Tzero.
  }
\end{equation}
This is the \textbf{pointwise surface SAR bound}, achieved for uniform
illumination of a flat surface or a perfectly focused beam.

\subsection{Lower bound via fixed patch}

For any specific patch area $A_0$,
\begin{equation}\label{eq:lower-bound}
  \psSAR \geq \frac{2\alpha}{\rhom}\;\langle\Sab\rangle_{A_0}\;
  h\!\Bigl(\frac{2\alpha\,m}{\rhom\,A_0}\Bigr).
\end{equation}
Choosing $A_0 = 4\;\mathrm{cm}^2$ (the ICNIRP averaging area) at 28\,GHz:
\begin{equation}\label{eq:lower-4cm2}
  \psSAR \geq 0.040 \times \max_{\vb{r}}\langle\Sab\rangle_{4\mathrm{cm}^2}(\vb{r}).
\end{equation}
This provides a computable lower bound from the ICNIRP spatial average.

\subsection{Uniform illumination on a flat surface}

For $\Sab(\vb{r}) = S_0$ everywhere on a flat surface, the area average
$\langle\Sab\rangle_A = S_0$ is independent of the patch. The optimization
reduces to $\max_{u > 0} h(u) = 1$ (achieved as $u \to 0^+$, i.e.\
$A \to \infty$), giving
\begin{equation}\label{eq:uniform-flat}
  \psSAR^{\mathrm{uniform}} = \frac{2\alpha}{\rhom}\,S_0
  = \SAR(\vb{r}, 0).
\end{equation}
The 10\,g averaging provides no reduction for uniform illumination.

\subsection{Thin skin depth limit}

As $\alpha \to \infty$ (skin depth $\to 0$), the SAR concentrates on the
surface: $\SAR(\vb{r},z) \to (\Sab(\vb{r})/\rhom)\,\delta(z)$ in the
distributional sense. The mass-averaged SAR over a slab of vanishing thickness
but finite area becomes
\begin{equation}\label{eq:thin-limit}
  \psSAR^{\delta\to 0}
  = \frac{1}{m}\;\sup_{\substack{A\;\text{connected}}}
    \int_A \Sab(\vb{r})\,\mathrm{d}A.
\end{equation}
When the entire illuminated surface fits within the 10\,g mass constraint,
this reduces to $\SAR_{\mathrm{wb}} = P_{\mathrm{abs}}/m$.

\subsection{The dilution factor}

\begin{definition}[Dilution factor]\label{def:dilution}
\begin{equation}\label{eq:dilution}
  \eta_{10\mathrm{g}}
  = \frac{\psSAR}{(2\alpha/\rhom)\,\max\Sab}
  = \sup_A\left[
    \frac{\langle\Sab\rangle_A}{\max\Sab}\;
    h\!\Bigl(\frac{2\alpha\,m}{\rhom\,A}\Bigr)
  \right].
\end{equation}
\end{definition}
By \cref{eq:upper-bound}, $\eta_{10\mathrm{g}} \leq 1$.
The dilution factor depends only on the \emph{shape} of the $\Sab$ distribution
(not its magnitude) and on the dimensionless ratio $2\alpha m/\rhom$.
Representative values:
\begin{itemize}[nosep]
\item Cosine distribution on a sphere ($R = 15$\,cm, 28\,GHz):
  $\eta_{10\mathrm{g}} \approx 0.67$.
\item Gaussian hotspot ($w = 10$\,mm) on a flat surface:
  $\eta_{10\mathrm{g}} \approx 0.03$.
\end{itemize}
The dilution factor decreases as the hotspot becomes more localized, because the
10\,g region must expand into lower-valued surface regions to satisfy the mass
constraint.

\subsection{Frequency dependence}

The prefactor $2\alpha/\rhom = 4\pi f\kappa/(c_0\rhom)$ grows with frequency
through the product $f\kappa$. Although $\kappa$ varies through the Cole-Cole
dispersion, the product $f\kappa$ increases monotonically above 6\,GHz for
biological tissue. This is one physical reason ICNIRP transitions from SAR-based
limits (below 6\,GHz) to $\Sab$-based limits (above 6\,GHz): the pointwise
surface SAR becomes very large even for modest $\Sab$, because absorption is
concentrated in an ever-thinner layer.


% ======================================================================
\section{Extension below 6~GHz}
\label{sec:below-6ghz}
% ======================================================================

The preceding analysis rests on the thin-skin approximation
$\alpha L \gg 1$, which permits replacing the exponential depth profile by
a surface delta function (\cref{eq:collapse}).  This approximation is
excellent above 6\,GHz, where the SAR penetration depth
$\delta_{\SAR} = 1/(2\alpha)$ is a fraction of a millimetre and all
transmitted power is absorbed within the top few per cent of the 10\,g
cube.  Below 6\,GHz, the penetration depth becomes comparable to or larger
than the cube side $L = 21.5\;\mathrm{mm}$, three new complications arise,
and the surface-only reduction fails.  This section quantifies the
transition, replaces the homogeneous half-space with a layered tissue
model, and derives the standing-wave SAR that results from internal
reflections.


% ----------------------------------------------------------------------
\subsection{Frequency transition}
\label{sec:freq-transition}
% ----------------------------------------------------------------------

\Cref{tab:skin-freq} collects the dielectric properties of skin tissue
from the Gabriel parametric model \cite{Gabriel1996} and the IT'IS tissue
database across the range 0.9--60\,GHz.  Three dimensionless quantities
govern the transition:

\begin{itemize}[nosep]
\item \textbf{Depth parameter} $2\alpha L$, comparing the SAR penetration
  depth to the cube side length.
\item \textbf{Depth efficiency} $h(2\alpha L)$, the fraction of the
  surface SAR that survives mass-averaging to depth~$L$.
\item \textbf{Capture fraction} $C = 1 - e^{-2\alpha L}$, the fraction of
  transmitted power absorbed within the cube depth.
\end{itemize}

\begin{table}[ht]
\centering
\caption{Skin tissue parameters across frequency.  Complex refractive
  index $\ntilde = n - j\kappa$, field attenuation rate
  $\alpha = k_0 \kappa$, SAR depth
  $\delta_{\SAR} = 1/(2\alpha)$, depth parameter $2\alpha L$ for the
  10\,g cube ($L = 21.5\;\mathrm{mm}$), depth efficiency $h(u)$, and
  capture fraction $C$.  Dielectric data: Gabriel~1996 /
  IT'IS.}
\label{tab:skin-freq}
\begin{tabular}{@{}r r r r r r r r r@{}}
\toprule
$f$ (GHz) & $\varepsilon_r$ & $\sigma$ (S/m) & $n$ & $\kappa$
  & $\alpha$ ($\mathrm{m^{-1}}$) & $2\alpha L$ & $h(2\alpha L)$ & $\delta_{\SAR}$ (mm) \\
\midrule
0.9  & 40.9 & 0.87  & 6.53 & 1.330 &   25  &  1.1 & 0.61 & 19.9 \\
2.45 & 38.0 & 1.46  & 6.22 & 0.860 &   44  &  1.9 & 0.45 & 11.3 \\
5    & 35.8 & 3.06  & 6.05 & 0.909 &   95  &  4.1 & 0.24 &  5.3 \\
10   & 31.3 & 8.01  & 5.73 & 1.256 &  263  & 11.3 & 0.088 &  1.9 \\
28   & 16.6 & 25.8  & 4.47 & 1.851 & 1086  & 46.7 & 0.021 &  0.46 \\
60   &  7.98 & 36.4 & 3.28 & 1.663 & 2092  & 89.9 & 0.011 &  0.24 \\
\bottomrule
\end{tabular}
\end{table}

\noindent
Three regimes emerge:

\paragraph{Thin skin ($\alpha L \gg 1$, above ${\sim}\,6\;\mathrm{GHz}$).}
The capture fraction $C \approx 1$: virtually all transmitted power is
absorbed within the cube depth.  The depth efficiency $h \ll 1$ and
$h(2\alpha L) \approx 1/(2\alpha L)$, so the mass-averaged SAR is simply
$\Sab/(\rhom L)$.  The energy conservation shortcut
(\cref{eq:SAR-energy}) applies directly.

\paragraph{Transition ($\alpha L \sim 1$, approximately
  1--6\,$\mathrm{GHz}$).}
The capture fraction drops below unity: at
$2.45\;\mathrm{GHz}$, $C = 0.85$, meaning 15\% of transmitted power
passes through the cube and is absorbed deeper in the body.  The depth
efficiency $h \approx 0.45$ is neither close to zero nor close to unity.
The mass-averaged SAR depends jointly on the surface distribution
\emph{and} the depth profile.  The thin-skin shortcut overpredicts
$\langle\SAR\rangle$ by $1/h \approx 2\times$; the half-space model
remains valid but the full depth integral must be retained.

\paragraph{Deep penetration ($\alpha L \ll 1$, below
  ${\sim}\,500\;\mathrm{MHz}$).}
The SAR varies slowly across the cube depth ($h \to 1$) and the
mass-averaged SAR approaches the surface SAR: $\langle\SAR\rangle \approx
(2\alpha/\rhom)\,\Sab$.  But now $\delta_{\SAR} \gg L$, so the
homogeneous half-space model itself breaks down: the field traverses
multiple tissue layers within the cube, and reflections from internal
interfaces (fat--muscle, bone--tissue) create standing waves that dominate
the SAR distribution.

\begin{remark}[Why the transition matters]
The transition regime 1--6\,GHz is precisely the frequency range where
ICNIRP~2020 switches from SAR-based to $\Sab$-based limits.  The
surface-only reduction of the previous sections provides
a clean analytical tool for the $\Sab$ regime.  Extending it below
6\,GHz requires accounting for finite penetration depth and internal
reflections, which is the subject of the remainder of this section.
\end{remark}


% ----------------------------------------------------------------------
\subsection{Layered tissue model}
\label{sec:layered}
% ----------------------------------------------------------------------

The homogeneous half-space treats tissue as a single medium with a single
attenuation constant.  Real tissue beneath the skin surface consists of
distinct layers with markedly different dielectric properties.  At
sub-6\,GHz frequencies the SAR penetration depth is comparable to or
larger than these layer thicknesses, so internal reflections are no longer
negligible.

We adopt a canonical three-layer planar model (normal incidence):

\begin{enumerate}[nosep]
\item \textbf{Skin} (epidermis + dermis): thickness
  $d_1 \approx 2\;\mathrm{mm}$, complex refractive index $\ntilde_1$,
  conductivity $\sigma_1$, mass density~$\rho_1$.
\item \textbf{Subcutaneous fat}: thickness
  $d_2 \approx 5$--$20\;\mathrm{mm}$ (varies by body region and
  individual), $\ntilde_2$, $\sigma_2$, $\rho_2$.
\item \textbf{Muscle} (semi-infinite): $\ntilde_3$, $\sigma_3$, $\rho_3$.
\end{enumerate}

\noindent
The dielectric contrast between these layers is extreme.  At
$2.45\;\mathrm{GHz}$:

\begin{center}
\begin{tabular}{@{}l c c c c c@{}}
\toprule
Tissue & $\varepsilon_r$ & $\sigma$ (S/m) & $n$ & $\kappa$ &
  $\delta_{\SAR}$ (mm) \\
\midrule
Skin   & 38.0 & 1.46 & 6.22 & 0.860 & 11.3 \\
Fat    &  5.3 & 0.08 & 2.31 & 0.127 & 76.5 \\
Muscle & 52.7 & 1.74 & 7.31 & 0.873 & 11.2 \\
\bottomrule
\end{tabular}
\end{center}

\noindent
Fat is nearly transparent ($\delta_{\SAR} = 76.5\;\mathrm{mm}$), while
skin and muscle are lossy.  This means the electromagnetic field can
traverse a thick fat layer with little attenuation and encounter the
fat--muscle interface, where the large impedance mismatch ($|r_{23}|
\approx 0.52$) generates a strong backward-travelling wave.

\paragraph{Generalized Fresnel reflection coefficients.}
Define the normal-incidence Fresnel reflection coefficient at the
interface between layers $j$ and $j+1$:
\begin{equation}\label{eq:rjj1}
  r_{j,\,j+1}
  = \frac{\ntilde_j - \ntilde_{j+1}}{\ntilde_j + \ntilde_{j+1}},
\end{equation}
and the one-way propagation factor across layer~$j$:
\begin{equation}\label{eq:Pj}
  P_j = e^{-jk_0\ntilde_j\,d_j}
  = e^{-\alpha_j d_j}\,e^{-j\beta_j d_j},
  \qquad
  |P_j| = e^{-\alpha_j d_j} \leq 1,
\end{equation}
where $\alpha_j = k_0\kappa_j$ and $\beta_j = k_0 n_j$.
The generalized (effective) reflection coefficient at each interface is
built recursively from the bottom up \cite{Chew1995}:

\begin{align}
  \widetilde{\Gamma}_3
    &= r_{2,3},
    \label{eq:Gamma3} \\[4pt]
  \widetilde{\Gamma}_2
    &= \frac{r_{1,2} + \widetilde{\Gamma}_3\,P_2^{\,2}}
            {1 + r_{1,2}\,\widetilde{\Gamma}_3\,P_2^{\,2}},
    \label{eq:Gamma2} \\[4pt]
  \widetilde{\Gamma}_1
    &= \frac{r_{0,1} + \widetilde{\Gamma}_2\,P_1^{\,2}}
            {1 + r_{0,1}\,\widetilde{\Gamma}_2\,P_1^{\,2}},
    \label{eq:Gamma1}
\end{align}

\noindent
where $r_{0,1} = (1 - \ntilde_1)/(1 + \ntilde_1)$ is the air-to-skin
coefficient and $P_j^{\,2} = e^{-2jk_0\ntilde_j d_j}$ is the round-trip
phase--attenuation factor across layer~$j$.  The overall power reflection
coefficient is $|\widetilde{\Gamma}_1|^2$, and the power transmission
coefficient replacing $\Tzero$ is
\begin{equation}\label{eq:T-layered}
  T_{\mathrm{lay}} = 1 - |\widetilde{\Gamma}_1|^2.
\end{equation}

\paragraph{Field in each layer.}
In layer~$j$ ($j = 1, 2, 3$), with $z$ measured as the local depth from
the top of that layer, the electric field takes the form
\begin{equation}\label{eq:Ej-layer}
  E_j(z) = A_j\bigl[e^{-jk_j z} + g_j\,e^{+jk_j z}\bigr],
\end{equation}
where $k_j = k_0\ntilde_j$ is the complex wave number in layer~$j$,
$A_j$ is the forward-wave amplitude at $z = 0$ (top of layer~$j$), and
\begin{equation}\label{eq:gj-def}
  \boxed{%
    g_j = \widetilde{\Gamma}_{j+1}\,P_j^{\,2}
    = \widetilde{\Gamma}_{j+1}\,e^{-2jk_j d_j}
  }
\end{equation}
is the \textbf{backward-to-forward amplitude ratio} at the top of
layer~$j$.  Physically, $g_j$ encodes the total reflected field returning
from all deeper interfaces (including multiple reflections), attenuated
and phase-shifted by the round trip across layer~$j$.

For muscle (layer~3, semi-infinite): $g_3 = 0$ because there is no
reflected wave returning from infinity.

\begin{remark}[Amplitude linkage between layers]
The forward amplitudes are related by the transmission coefficient and
propagation factor at each interface.  Continuity of the tangential
electric field at the top of layer~$j+1$ (bottom of layer~$j$) gives
\begin{equation}\label{eq:Aj-recursion}
  A_{j+1}(1 + g_{j+1})
  = A_j\bigl[e^{-jk_j d_j} + g_j\,e^{+jk_j d_j}\bigr]
  = A_j\,e^{-jk_j d_j}\bigl[1 + g_j\,e^{2jk_j d_j}\bigr],
\end{equation}
from which $A_{j+1}$ follows once $A_1$ is known from the incident field
and the overall transmission through the air--skin interface.
\end{remark}


% ----------------------------------------------------------------------
\subsection{Standing-wave SAR}
\label{sec:standing-wave}
% ----------------------------------------------------------------------

The field intensity in layer~$j$ follows from \cref{eq:Ej-layer}:
\begin{equation}\label{eq:E2-layer}
  |E_j(z)|^2
  = |A_j|^2\bigl|e^{-jk_j z} + g_j\,e^{+jk_j z}\bigr|^2.
\end{equation}
Expanding the squared modulus (with $k_j = \beta_j - j\alpha_j$ in the
engineering $e^{+j\omega t}$ convention), the three terms are:

\begin{align}
  \bigl|e^{-jk_j z} + g_j\,e^{+jk_j z}\bigr|^2
  &= e^{-2\alpha_j z}
   + |g_j|^2\,e^{+2\alpha_j z}
   + 2|g_j|\cos(2\beta_j z + \psi_j),
   \label{eq:E2-three-terms}
\end{align}

\noindent
where $\psi_j = \arg(g_j)$.  The SAR in layer~$j$ is therefore

\begin{equation}\label{eq:SAR-standing}
  \boxed{%
  \SAR_j(z)
  = \frac{\sigma_j\,|A_j|^2}{2\rho_j}
    \Bigl[
      \underbrace{e^{-2\alpha_j z}}_{\text{forward}}
    + \underbrace{|g_j|^2\,e^{+2\alpha_j z}}_{\text{backward}}
    + \underbrace{2|g_j|\cos(2\beta_j z + \psi_j)}_{\text{interference}}
    \Bigr].
  }
\end{equation}

\noindent
The three terms have distinct physical origins:

\begin{enumerate}[nosep]
\item \textbf{Forward decay.}  The incident wave attenuating
  exponentially into the layer, as in the homogeneous half-space model.
  This is the only term when $g_j = 0$ (no reflections), recovering
  the single-exponential SAR of \cref{eq:SAR-depth}.

\item \textbf{Backward growth.}  The wave reflected from deeper
  interfaces, growing exponentially toward the surface (it has
  already been attenuated on its inward journey and is now propagating
  outward through layer~$j$).  The factor $|g_j|^2$ is the round-trip
  power attenuation of the reflected wave.

\item \textbf{Interference.}  The coherent beating between the forward
  and backward waves, producing a standing-wave pattern with spatial
  period $\pi/\beta_j = \lambda_0/(2n_j)$.  At $2.45\;\mathrm{GHz}$ in
  skin, this period is $\pi/319.6 \approx 9.8\;\mathrm{mm}$.  The
  interference oscillation can shift the SAR maximum away from the
  surface if the backward wave is sufficiently strong.
\end{enumerate}

\begin{remark}[Recovery of the half-space limit]
When there are no internal reflections ($g_j = 0$ for all~$j$),
\cref{eq:SAR-standing} reduces to $\SAR = (\sigma|A|^2/2\rho)\,
e^{-2\alpha z}$, which is \cref{eq:SAR-depth} with the identification
$|A|^2 = (4\alpha/\sigma)\,\Sab$.  The layered model is therefore a
strict generalization of the half-space theory.
\end{remark}


% ----------------------------------------------------------------------
\subsection{Integrated SAR over the cube}
\label{sec:layered-integral}
% ----------------------------------------------------------------------

The 10\,g cube extends from $z = 0$ (the body surface) to $z = L$.  When
$L$ exceeds $d_1$ or $d_1 + d_2$, the cube spans multiple tissue layers.
The absorbed energy in each layer segment is computed by integrating
\cref{eq:SAR-standing} analytically.

\paragraph{Finite layer ($g_j \neq 0$, thickness $d_j$).}
The integral of $\SAR_j \rho_j$ over the full extent of layer~$j$ is

\begin{equation}\label{eq:layer-integral}
  \boxed{%
  \int_0^{d_j}\!\SAR_j\,\rho_j\,\mathrm{d}z
  = \frac{\sigma_j |A_j|^2}{2}
    \biggl[
      \frac{1 - e^{-2\alpha_j d_j}}{2\alpha_j}
    + |g_j|^2\,\frac{e^{2\alpha_j d_j} - 1}{2\alpha_j}
    + \frac{|g_j|}{\beta_j}\Bigl(
        \sin(2\beta_j d_j + \psi_j) - \sin\psi_j
      \Bigr)
    \biggr].
  }
\end{equation}

\noindent
The three terms correspond to the forward, backward, and interference
contributions respectively.  Each has a clear physical interpretation:

\begin{itemize}[nosep]
\item The forward term $\propto (1 - e^{-2\alpha_j d_j})/(2\alpha_j)$ is
  the energy absorbed from the incident wave.  For thick layers
  ($\alpha_j d_j \gg 1$) it saturates at $1/(2\alpha_j)$.
\item The backward term $\propto |g_j|^2(e^{2\alpha_j d_j} - 1)/(2\alpha_j)$
  is the energy deposited by the reflected wave as it propagates upward
  through the layer.  It can be large when $|g_j|$ is significant and the
  layer is thick.
\item The interference term oscillates with the phase traversal
  $2\beta_j d_j$ and can be positive or negative, adding or subtracting
  from the total depending on the standing-wave phase.
\end{itemize}

\paragraph{Semi-infinite layer ($g_3 = 0$, remaining depth $d_3$).}
For the muscle layer, which has no reflected wave, the integral from the
top of layer~3 over the remaining cube depth
$d_3 = L - d_1 - d_2$ simplifies to

\begin{equation}\label{eq:muscle-integral}
  \int_0^{d_3}\!\SAR_3\,\rho_3\,\mathrm{d}z
  = \frac{\sigma_3\,|A_3|^2}{2}\;\frac{1 - e^{-2\alpha_3 d_3}}{2\alpha_3}.
\end{equation}

\paragraph{Total mass-averaged SAR.}
The peak spatial-average SAR over the cube is the sum of the layer
contributions divided by the total tissue mass:
\begin{equation}\label{eq:psSAR-layered}
  \langle\SAR\rangle_{\mathrm{cube}}
  = \frac{1}{m}\sum_{j=1}^{3}
    \int_0^{\min(d_j,\,z_j^{\max})}\!\SAR_j\,\rho_j\,\mathrm{d}z,
\end{equation}
where $z_j^{\max}$ is the remaining cube depth available for layer~$j$
after subtracting the thicknesses of the layers above.  If the layer
densities are approximately equal ($\rho_1 \approx \rho_2 \approx \rho_3
\approx \rhom$), then $m \approx \rhom L^3$ and the formula simplifies.
In practice, fat density ($\rho_2 \approx 920\;\mathrm{kg/m^3}$) differs
from skin and muscle ($\rho_1, \rho_3 \approx 1100\;\mathrm{kg/m^3}$) by
roughly 15\%, which introduces a corresponding correction to the mass
normalization.

\begin{remark}[Comparison with the homogeneous result]
In the half-space model, the analogous expression is simply
$(2\alpha/\rhom)\,h(2\alpha L)\,\Sab$.  The layered formula
\cref{eq:psSAR-layered} replaces this single term with a sum of
three-term integrals per layer, each depending on the local tissue
properties and the backward-to-forward ratio $g_j$.  The additional
complexity is the price of accounting for internal reflections.
\end{remark}


% ----------------------------------------------------------------------
\subsection{Subsurface SAR peaks}
\label{sec:subsurface-peaks}
% ----------------------------------------------------------------------

In the homogeneous half-space model, the SAR maximum always occurs at the
body surface ($z = 0$), because the exponential $e^{-2\alpha z}$ is
monotonically decreasing.  With internal reflections this is no longer
guaranteed.  The backward-travelling wave can constructively interfere
with the forward wave at some depth inside the tissue, creating a SAR
maximum below the surface.

\paragraph{Criterion for a subsurface peak in layer~1.}
Differentiating \cref{eq:SAR-standing} in layer~1:
\begin{equation}\label{eq:dSAR-dz}
  \frac{\mathrm{d}\SAR_1}{\mathrm{d}z}\bigg|_{z=0}
  = \frac{\sigma_1 |A_1|^2}{2\rho_1}
    \bigl[
      -2\alpha_1
      + 2\alpha_1 |g_1|^2
      - 4\beta_1 |g_1|\sin\psi_1
    \bigr].
\end{equation}
For the SAR to increase with depth at the surface, the derivative must be
positive.  Dividing by $2\alpha_1 > 0$:

\begin{equation}\label{eq:subsurface-condition}
  \boxed{%
    |g_1|^2 - \frac{2\beta_1}{\alpha_1}\,|g_1|\sin\psi_1 > 1
    \qquad\Longleftrightarrow\qquad
    \text{SAR maximum is below the surface.}
  }
\end{equation}

\noindent
This condition involves two factors:

\begin{itemize}[nosep]
\item \textbf{Reflection strength} ($|g_1|^2$).  For this term alone to
  exceed unity, the backward wave at the surface must carry more power
  than the forward wave, which is uncommon.  Typical values at
  $2.45\;\mathrm{GHz}$ are $|g_1| \approx 0.6$, so $|g_1|^2 \approx 0.37
  < 1$.

\item \textbf{Interference sign} ($-\sin\psi_1$).  When
  $\psi_1 = \arg(g_1)$ is negative (which is typical, since the
  propagation phase $-2k_1 d_1$ rotates $\arg(g_1)$ into the lower half
  of the complex plane), $\sin\psi_1 < 0$ and the interference term
  $-(2\beta_1/\alpha_1)|g_1|\sin\psi_1$ becomes positive.  Because
  $\beta_1/\alpha_1 = n_1/\kappa_1 \approx 7$ for skin at
  $2.45\;\mathrm{GHz}$, even a modest reflection $|g_1| \approx 0.6$ can
  produce a large interference contribution.
\end{itemize}

\paragraph{Numerical evaluation at $2.45\;\mathrm{GHz}$.}
For a skin--fat--muscle stack with $d_1 = 2\;\mathrm{mm}$,
$d_2 = 10\;\mathrm{mm}$, the computed backward-to-forward ratio is
$|g_1| \approx 0.61$, $\psi_1 \approx -61^\circ$.  The subsurface
condition evaluates to
\[
  |g_1|^2 - \frac{2\beta_1}{\alpha_1}\,|g_1|\sin\psi_1
  \;=\; 0.37 + 7.7
  \;=\; 8.0 \;\gg\; 1.
\]
The condition is satisfied by a wide margin, confirming that the SAR
maximum at $2.45\;\mathrm{GHz}$ can lie below the skin surface.

\paragraph{The role of the fat layer.}
The fat layer plays a distinctive role in creating subsurface SAR
peaks.  Its conductivity is very low ($\sigma_{\mathrm{fat}} \approx
0.08\;\mathrm{S/m}$ at $2.45\;\mathrm{GHz}$, compared to
${\sim}\,1.5\;\mathrm{S/m}$ for skin and muscle), so the SAR within the
fat itself is small.  But the field amplitude can remain high throughout
the fat because of the low attenuation ($\alpha_{\mathrm{fat}} d_2
\approx 0.07$ for a $10\;\mathrm{mm}$ layer, corresponding to only
13\% power loss across the layer).

\begin{theorem}[Subsurface peak mechanism]\label{thm:subsurface}
At sub-GHz to low-GHz frequencies with thin skin over thick
subcutaneous fat over muscle, the peak SAR within the 10\,g cube can
exceed the surface SAR.  Two effects cooperate:
\begin{enumerate}[nosep]
\item The field partially traverses the low-loss fat layer without
  substantial attenuation.
\item The fat--muscle interface reflection ($|r_{2,3}| \approx 0.52$ at
  $2.45\;\mathrm{GHz}$) creates a backward wave that interferes
  constructively with the forward wave near the skin--fat boundary.
\end{enumerate}
The constructive interference produces SAR enhancements of order
$(1 + |g_1|)^2/(1 - |g_1|)^2$ at the standing-wave antinodes, which can
reach factors of 2--4 relative to the pure-decay prediction.
\end{theorem}

\begin{remark}[Implication for the surface-based reduction]
The existence of subsurface SAR peaks means that the surface quantity
$\Sab$ alone does not bound the volumetric SAR.  The reduction
$\psSAR = f(\Sab)$ of \cref{thm:main,thm:cube} requires the
homogeneous half-space assumption ($g = 0$, monotonic depth decay).
Below 6\,GHz, the layered model must be used, and the ``input'' to the
SAR computation is no longer $\Sab$ on the body surface but rather the
full set of forward amplitudes $\{A_j\}$ and backward-to-forward ratios
$\{g_j\}$ determined by the tissue stratigraphy.
\end{remark}


% ----------------------------------------------------------------------
\subsection{When the theory breaks}
\label{sec:limitations-low-freq}
% ----------------------------------------------------------------------

The layered half-space model of the preceding subsections extends the
analysis to the low-GHz range but still relies on a locally planar,
stratified geometry.  Three phenomena invalidate this picture at
sufficiently low frequencies.

\paragraph{Whole-body resonance (below ${\sim}\,300\;\mathrm{MHz}$).}
When the free-space wavelength becomes comparable to the body height
($\lambda/2 \approx 0.85\;\mathrm{m}$ at ${\sim}\,175\;\mathrm{MHz}$
for an adult, resonance peaked near $70\;\mathrm{MHz}$
for a grounded standing adult \cite{Durney1986}), the human body acts as
a resonant dipole antenna.  Internal fields are dominated by whole-body
current distributions, not by the surface-transmitted plane wave.
Localized hotspots form at anatomical constrictions (ankles, wrists, neck)
where current density is highest, and these hotspot locations bear no
relation to the surface absorbed power density.  The surface-to-volume
reduction is fundamentally inapplicable in this regime.

\paragraph{Finite-thickness structures.}
Thin anatomical features such as the ear pinna (cartilage thickness
${\sim}\,2\;\mathrm{mm}$) violate the semi-infinite half-space
assumption even at higher frequencies.  At $2\;\mathrm{GHz}$,
$\alpha d \approx 0.09$ for a $2\;\mathrm{mm}$ ear, so the one-way field
attenuation is only $e^{-\alpha d} \approx 0.91$: the transmitted field
partially exits the far side of the ear and illuminates the opposite
surface.  The SAR in such structures requires a finite-slab transmission
model rather than a half-space model.

\paragraph{Multi-path and deep hotspots.}
At frequencies where the penetration depth exceeds the body diameter at
narrow cross-sections (e.g., $\delta_{\SAR} \approx 20\;\mathrm{mm}$ at
$900\;\mathrm{MHz}$ versus forearm diameter ${\sim}\,60\;\mathrm{mm}$),
the field transmitted through one side of the body can reach the opposite
side.  Reflections from the far body--air interface create a second
backward wave, and the resulting interference pattern can produce deep
hotspots at locations not predictable from the surface fields alone.  This
``multi-path through the body'' effect is distinct from the layered-tissue
standing waves discussed above: it depends on the full 3D body geometry,
not just the local tissue stratigraphy.




% ======================================================================
\section{Summary}
\label{sec:summary}
% ======================================================================

Two paths lead from $\Sab$ on the body surface to \psSAR:

\paragraph{Path A (exact depth integral).}
\begin{equation}\label{eq:chain}
  \Sab(\vb{r})
  \;\xrightarrow{\text{\cref{eq:E-depth}}}\;
  |\vb{E}(\vb{r},z)|^2 = \frac{4\alpha}{\sigma}\Sab\,e^{-2\alpha z}
  \;\xrightarrow{\text{\cref{eq:SAR-depth}}}\;
  \SAR(\vb{r},z) = \frac{2\alpha}{\rhom}\Sab\,e^{-2\alpha z}
  \;\xrightarrow{\text{\cref{eq:main}}}\;
  \psSAR.
\end{equation}
This yields the slab result (\cref{thm:main}) with the depth efficiency $h(u)$.
It is exact under the universal-$\alpha$ and half-space approximations.

\paragraph{Path B (energy conservation).}
In the thin-skin limit ($\alpha L \gg 1$, always satisfied above 6\,GHz), the
3D integral collapses to a surface integral
(\cref{eq:collapse}): absorbed power in the cube equals $\Sinc\Tzero$ times
the cube's projected area (\cref{eq:SAR-energy}).  This bypasses the
depth coordinate entirely and gives a closed-form result for the standard
cube that depends only on $\khat$ and $L$.

\medskip\noindent
Both paths rest on three approximations:
\begin{enumerate}[nosep]
\item the universal $\alpha$ approximation (angle-independent decay rate),
  error $\leq 2.2$\% at 28\,GHz;
\item the flat-slab volume approximation (neglecting curvature corrections to
  mass), error $\mathcal{O}(d/R)$;
\item the Fresnel half-space model (neglecting subsurface reflections), valid
  whenever skin depth $\ll$ body dimensions.
\end{enumerate}
Path~B additionally requires $\alpha L \gg 1$ and a body large enough to fill
the cube's cross-section.  For thin or folded anatomy (ears, fingers), the
general surface integral (\cref{eq:SAR-general}) is needed.

The key numerical quantities at 28\,GHz for skin tissue are collected in
\cref{tab:summary}.

\begin{table}[ht]
\centering
\caption{Summary of key quantities at 28\,GHz (skin, $\rhom = 1000\;\mathrm{kg/m^3}$).}
\label{tab:summary}
\begin{tabular}{@{}lll@{}}
\toprule
Quantity & Symbol & Value \\
\midrule
Complex refractive index & $\ntilde$ & $4.49 - 1.79j$ \\
Normal-incidence transmission & $\Tzero$ & 0.539 \\
Field decay rate & $\alpha = k_0\kappa$ & $1047\;\mathrm{m^{-1}}$ \\
Power $e$-folding depth & $1/(2\alpha)$ & 0.48\,mm \\
SAR-to-APD prefactor & $2\alpha/\rhom$ & 2.09 \\
\addlinespace
Cube side length & $L = (m/\rhom)^{1/3}$ & 21.5\,mm \\
Cube depth parameter & $u_{\mathrm{cube}} = 2\alpha L$ & 45.1 \\
\addlinespace
Cube SAR (axis-aligned, energy cons.) & $\Sinc\Tzero/(\rhom L)$ & $0.047\,\Sinc\Tzero$ \\
\addlinespace
$\psSAR^{\mathrm{cube}}$ for $\Sab = 20\;\mathrm{W/m^2}$ (uniform) & & $0.93\;\mathrm{W/kg}$ \\
$\psSAR$ upper bound for $\Sab = 20\;\mathrm{W/m^2}$ & & $41.9\;\mathrm{W/kg}$ \\
Slab, optimal cap on $R=15$\,cm sphere, $\Sinc=10\;\mathrm{W/m^2}$ & & $7.59\;\mathrm{W/kg}$ at $\theta_0 \approx 49^\circ$ \\
\bottomrule
\end{tabular}
\end{table}


% ======================================================================
% References
% ======================================================================

\begin{thebibliography}{9}

\bibitem{ICNIRP2020}
ICNIRP,
``Guidelines for limiting exposure to electromagnetic fields
(100\,kHz to 300\,GHz),''
\emph{Health Physics}, vol.~118, no.~5, pp.~483--524, 2020.

\bibitem{ICNIRPdiff}
ICNIRP,
``Differences between the ICNIRP (2020) and previous guidelines,''
\url{https://www.icnirp.org/en/differences.html}, accessed April 2026.
%
\emph{Key quotes:}
``ICNIRP (2020) requires [SAR] to be averaged over a 10-g cubic region''
(changed from ``contiguous tissue region'' in 1998);
``This 4-cm$^2$ averaging area matches the face of the averaging volume
(10\,g) of SAR.''

\bibitem{C953}
IEEE,
``IEEE~C95.3-2021 -- Recommended practice for measurements and computations
of electric, magnetic, and electromagnetic fields with respect to human
exposure to such fields, 0\,Hz to 300\,GHz,'' 2021.
%
Annex~E specifies the averaging volume construction algorithm.

\bibitem{62704-1}
IEC/IEEE~62704-1:2017,
``Determining the peak spatial-average specific absorption rate (SAR) in the
human body from wireless communications devices, 30\,MHz to 6\,GHz --
Part~1: General requirements for using the finite-difference time-domain
(FDTD) method for SAR calculations,'' 2017.
%
The model space is subdivided into voxels aligned with the global coordinate
system; cubes are expanded evenly in six axis-aligned directions from each
seed voxel until the enclosed tissue mass is within $0.0001$\% of the
prescribed mass.

\bibitem{62704-4}
IEC/IEEE~62704-4:2020,
``Determining the peak spatial-average specific absorption rate (SAR) in the
human body from wireless communication devices, 30\,MHz to 6\,GHz --
Part~4: General requirements for using the finite element method for SAR
calculations,'' 2020.

\bibitem{62209-1528}
IEC/IEEE~62209-1528:2020,
``Measurement procedure for the assessment of specific absorption rate of
human exposure to radio frequency fields from hand-held and body-worn
wireless communication devices -- Part~1528: Human models, instrumentation
and procedures (4\,MHz to 10\,GHz),'' 2020.

\bibitem{Gabriel1996}
C.~Gabriel, S.~Gabriel, and E.~Corthout,
``The dielectric properties of biological tissues: I.~Literature survey,''
\emph{Phys.\ Med.\ Biol.}, vol.~41, no.~11, pp.~2231--2249, 1996.

\bibitem{Chew1995}
W.~C.~Chew,
\emph{Waves and Fields in Inhomogeneous Media},
IEEE Press, 1995.

\bibitem{Durney1986}
C.~H.~Durney, H.~Massoudi, and M.~F.~Iskander,
\emph{Radiofrequency Radiation Dosimetry Handbook},
USAF School of Aerospace Medicine, 4th~ed., 1986.

\end{thebibliography}

\end{document}
